\chapter{Stato dell'arte e Background}
\label{capitolo2}
 

%\noindent Nella seconda sezione si riporta lo stato dell'arte del settore, un inquadramento dell'area di ricerca orientato a portare il lettore all'interno della problematica affrontata. Bisogna dimostrare di conoscere le cose fatte fino ad ora in questo campo e il perch\'e si sia reso necessario lo svolgimento di questo lavoro. Questa sezione deve essere grondante di citazioni bibliografiche~\cite{marco02bud}.

\section{Stato dell'arte}
All'interno delle reti sociali sta sempre pi\'u prendendo piede il problema della polarizzazione delle opinioni. Nel linguaggio comune il confronto tra individui ha sempre generato una forte controversia nelle opinioni oppure una visione comune. 
I social network hanno permesso all'utente di poter diffondere attraverso post, messaggi o espressioni audio video le proprie opinioni e/o pensieri all'interno di una comunità sociale. A tal proposito per favorire la diffusione delle diverse correnti di pensiero i social network stanno sempre più sviluppando algoritmi per permettere di identificare le comunità isolate che condividono un unico punto di vista di un problema.
La polarizzazione è un algoritmo matematico che applicato all'interno delle rete sociali consente di capire quanto un utente che accede per la prima volta all'interno di una rete sociale, venga influenzato dagli altri utenti e quanto una news o un giudizio si propaga all'interno di una rete sociale.
Prima di poter illustrare questo algoritmo con le relative problematiche verrà illustrata una definizione di rete sociale.

\subsection{Sociale}
Una rete sociale consiste in un qualsiasi gruppo di individui connessi tra loro da diversi legami sociali. Per gli esseri umani i legami vanno dalla conoscenza casuale, ai rapporti di lavoro, ai vincoli familiari. Le reti sociali sono spesso usate come base di studi interculturali in sociologia, in antropologia, in etologia.

L'analisi delle reti sociali, ovvero la mappatura e la misurazione delle reti sociali, può essere condotta con un formalismo matematico usando la teoria dei grafi. In generale, il corpus teorico ed i modelli usati per lo studio delle reti sociali sono compresi nella cosiddetta \textit{social network analysis}.

La ricerca condotta nell'ambito di diversi approcci disciplinari ha evidenziato come le reti sociali operino a pi\'u livelli e svolgano un ruolo cruciale nel determinare le modalità di risoluzione di problemi e i sistemi di gestione delle organizzazioni, nonché le possibilità dei singoli individui di raggiungere i propri obiettivi.

\subsection{Le reti}
La diffusione del web e del termine social network ha creato negli ultimi anni alcune ambiguità di significato. La rete sociale è infatti storicamente, in primo luogo, una rete fisica.

Rete sociale è, ad esempio, una comunità di lavoratori, che si incontra nei relativi circoli dopolavoristici e che costituisce una delle associazioni di promozione sociale. Esempi di reti sociali sono inoltre le comunità di sportivi, attivi o sostenitori di eventi, le comunità unite da problematiche strettamente lavorative e di tutela sindacale del diritto nel lavoro, le confraternite e in generale le comunità basate sulla pratica comune di una religione e il ritrovo in chiese, templi, moschee, sinagoghe e altri luoghi di culto.

La rete sociale ha una sua versione all'interno di \textbf{Internet} definita anche come \textit{Social media}, questa è una delle forme più evolute di comunicazione in rete. La rete delle relazioni sociali che ciascuno di noi tesse ogni giorno, in maniera più o meno casuale, nei vari ambiti della nostra vita, si può così "materializzare", organizzare in una "mappa" consultabile, e arricchire di nuovi contatti. I principali social network sono: \textit{Facebook, MySpace, Instagram, Twitter, Google+, LinkedIn, Ask.fm, Pinterest, Formspring, Bebo, Friendster, Hi5, Ning, Tagged, Meetup, Tumblr}.

\subsection{Twitter}
\textit{Twitter} è il social network di riferimento utilizzato per la realizzazione di questa tesi di laurea, ne diamo una breve descrizione, è una rete sociale, creata il 21 marzo 2006 dalla \textit{Obvious Corporation} di San Francisco, che fornisce agli utenti, attraverso l'omonima piattaforma, una pagina personale aggiornabile tramite messaggi di testo con lunghezza massima di 140 caratteri; nel 2017 l'azienda ha aumentato la lunghezza dei tweet a 280 caratteri per alcuni paesi. Gli aggiornamenti di stato possono essere effettuati tramite il sito stesso, via SMS, con programmi di messaggistica istantanea, posta elettronica, oppure tramite varie applicazioni basate sulle API di Twitter. 

Il nome "Twitter" deriva dal verbo inglese to tweet che significa "cinguettare". Tweet è anche il termine tecnico degli aggiornamenti del servizio. I Tweet che contengono esattamente 140 caratteri vengono chiamati \textit{Twoosh}. Gli aggiornamenti sono mostrati nella pagina di profilo dell'utente e comunicati agli utenti che si sono registrati per riceverli. È anche possibile limitare la visibilità dei propri messaggi oppure renderli visibili a chiunque. \begin{figure}[h!]
    \begin{center}
      \includegraphics[scale= 0.15]{twitter-logo.png}
	%\psfig{file=./pictures/logo}
    \end{center}
  \end{figure}

Questo social network è molto utilizzato da personaggi famosi, sportivi, ma anche capi di stato, esponenti politici ed economici, perfino il Papa, consentendo a tutti gli utenti iscritti alla piattaforma di poter visualizzare le idee ed opinioni espresse da questi personaggi. Proprio per questo motivo \textit{Twitter} consente una diffusione su larga scala di idee ed opinioni da parte di moltissimi utenti, gli stessi giornalisti utilizzano questa piattaforma per diffondere \textit{scoop} o notizie di vario genere, proprio per la grande visibilità che fornisce questo social network.
Vengono generati moltissimi dati ed è proprio per questo motivo che ho deciso di utilizzarlo strumento per la raccolta dati e cercare di calcolare la polarizzazione, permettendomi di visualizzare come le informazioni espressi si propagano all'interno della rete e come queste influenzano il pensiero degli utenti.

\subsubsection{Tweet}
L'elemento alla base di questo social network è il \textit{Tweet}, che non è altro che un messaggio contenente le opinioni, le idee, i pensieri dell'utente. Tutte queste cose possono essere arricchite utilizzato elementi multimediali come video, immagini oppure link a pagine web. Il tweet una volta pubblicato potrà essere accessibile per tutti gli utenti delle rete che hanno una relazione di "amicizia" con l'utente in questione, potendo anche effettuare delle operazioni aggiuntive quale commentare la notizia, apprezzarla oppure \textit{retwettarla} cioè condividere queste informazioni a tutti i propri collegamenti.
I Tweet possono essere anche etichettati utilizzando due strumenti messi a disposizione da questo social network:
\begin{itemize}
\item \textit{Mentions}: sono dei riferimento ad utenti, cioè come assegnare un mittente ad un email, affermando all'interno della rete che quel tweet è destinato esclusivamente a quella persona indicata attraverso questo strumento. Può essere utile per condividere più informazioni con gli amici delle rete più stretti, oppure per formulare degli attacchi, riflessioni o considerazioni anche con utenti famosi. Per fare questa operazione è necessario inserire la \@ ed il nome dell'utente in questione per poter collegare al tweet in questione alla persona desiderata.
\item \textit{Hashtag}: parole o combinazioni di parole concatenate precedute dal simbolo cancelletto (\#).
Etichettando un messaggio con un hashtag si crea un collegamento ipertestuale a tutti i messaggi recenti che citano lo stesso hashtag. Molto  vengono utilizzati durante \textit{Show Tv}, \textit{comizi}, \textit{pubblicità}, \textit{eventi politici} per poter consentire alle persone che vogliano esprimere la loro opinioni sugli argomenti espressi durante questi eventi circoscrivendo i Tweet. 

\end{itemize}

\section{Polarizzazione}
